%% LyX 2.4.4 created this file.  For more info, see https://www.lyx.org/.
%% Do not edit unless you really know what you are doing.
\documentclass[oneside,english]{amsart}
\usepackage[T1]{fontenc}
\usepackage[latin9]{inputenc}
\usepackage{enumitem}
\usepackage{amsthm}

\makeatletter
%%%%%%%%%%%%%%%%%%%%%%%%%%%%%% Textclass specific LaTeX commands.
\numberwithin{equation}{section}
\numberwithin{figure}{section}
\newlength{\lyxlabelwidth}      % auxiliary length 

%%%%%%%%%%%%%%%%%%%%%%%%%%%%%% User specified LaTeX commands.
\usepackage{enumitem}
\usepackage{ifthen}
\usepackage{lastpage}
%sets math fonts to be more readable
\everymath=\expandafter{\the\everymath\displaystyle}
%\DeclareMathSizes{10}{10}{10}{10}
\usepackage{multicol}
\usepackage{fancyhdr}
\pagestyle{fancy}
\usepackage{titlepic}
\usepackage{graphicx}
\usepackage{lineno}
\usepackage{hyperref}
%allows header and footer on titlepage by redefining the plain pagestyle issued by \maketitle to be fancy
%\makeatletter
%\let\ps@plain\ps@fancy
%\makeatother
\usepackage{advdate}    % Advancing/saving dates
\usepackage{datetime}   % Dates formatting
\usepackage{datenumber} % Counters for dates

\usepackage{hyperref}
\usepackage[usenames,dvipsnames]{xcolor} %adds additional color options
\hypersetup{colorlinks=true, urlcolor=Blue, linkcolor=Blue, citecolor=Blue}
\usepackage{url}
% Added by lyx2lyx
\setlength{\parskip}{\medskipamount}
\setlength{\parindent}{0pt}

\makeatother

\usepackage{babel}
\begin{document}
\renewcommand{\author}{Dr.\,James Ashe}

\newcommand{\classNum}{137}

\newcommand{\classSec}{all sections}

\newcommand{\nameType}{Math Lab Project}

\newcommand{\examType}{}

\renewcommand{\date}{}

%%First page's header%%%%%%%%%%%%%%%%%%%%%%
\fancypagestyle{firststyle} %%header settings for first pagr
{
	\fancyhead[L]{MTH \classNum{}(\classSec{}) \author{} \textbf{\examType{}} \date{}}
	\fancyhead[C]{}
	\fancyhead[R]{\textbf{\nameType}}
	\setlength{\headsep}{5pt}
}
%%%%%%%%%%%%%%%%%%%%%%%%%%%%%%%%%%%%%%%%%%

%%Next page's header%%%%%%%%%%%%%%%%%%%%%%%
\setlist{nolistsep}
\lhead{\classNum{}(\classSec{})} 
\chead{\textbf{\nameType}} 
\cfoot{\thepage{}~of~\pageref{LastPage}} 
\setlength{\headsep}{5pt} 
%%%%%%%%%%%%%%%%%%%%%%%%%%%%%%%%%%%%%%%%%%%

%%Various settings; see the preamble for other settings%%%%%
%%\setenumerate[0]{leftmargin=12pt,itemindent=0pt,labelwidth=0pt}
\setlist{itemsep=2pt}
\renewcommand{\labelenumi}{\textbf{\arabic{enumi}.}}
\renewcommand{\labelenumii}{\textbf{\alph{enumii}.}}
\renewcommand{\labelenumiii}{\textbf{\roman{enumiii}.}}
\thispagestyle{firststyle}
%%%%%%%%%%%%%%%%%%%%%%%%%%%%%%%%%%%%%%%%%%

\*

\section*{Quadratic Function Project}
\begin{enumerate}[resume]
\item Find a parabola in a work of art or architecture. Take or find a picture
that gives a good frontal view of the parabola. (Here is a website
you may want to take a look at: \hyperref[http://compfight.com/search/parabola/1-2-1-1]{http://compfight.com/search/parabola/1-2-1-1}.)\vfill
\item Draw a coordinate system over the picture (you may need to enlarge
the photo; here is a free photo editor \hyperref[http://www.gimp.org/downloads/]{http://www.gimp.org/downloads/}).

\begin{enumerate}
\item Mark the \emph{x }and \emph{y }axes. 
\item Mark the scale.\vfill
\end{enumerate}
\item Find three points on the graph that can be used to find the equation
of the parabola in the form $f(x)=a\left(x-h\right)^{2}+k$.\vfill
\item Use three different points on the graph to check your equation (they
only need to be approximations). \vfill
\item Identify the following:

\begin{enumerate}
\item The vertex.
\item All intercepts.
\item The maximum and minimum.
\item The domain and the range.\vfill
\end{enumerate}
\item Write the equation in the form $f(x)=ax^{2}+bx+c$. Show all your
work.\vfill
\item Write an essay or prepare a presentation. 

\begin{enumerate}
\item Describe how you completed each of the above steps.
\item Show all your work from the above steps.
\item Include the picture with the coordinate system.\vfill
\end{enumerate}
\end{enumerate}

\end{document}
